\documentclass[11pt,reqno]{amsart}
\usepackage[utf8]{inputenc}
\usepackage{amsmath,amssymb,amsthm,mathtools}
\usepackage[margin=1in]{geometry}
\usepackage[colorlinks=true,linkcolor=blue,citecolor=blue,urlcolor=blue]{hyperref}

\theoremstyle{plain}
\newtheorem{theorem}{Theorem}
\newtheorem{proposition}[theorem]{Proposition}
\newtheorem{lemma}[theorem]{Lemma}
\newtheorem{corollary}[theorem]{Corollary}
\newtheorem{conjecture}[theorem]{Conjecture}
\theoremstyle{definition}
\newtheorem{definition}[theorem]{Definition}
\newtheorem{remark}[theorem]{Remark}

% Named, unnumbered statements, to preserve the paper's A/B/C and Lemma 1/2 naming.
\theoremstyle{plain}
\newtheorem*{thmA}{Theorem A (slab identity)}
\newtheorem*{thmB}{Theorem B}
\newtheorem*{thmC}{Theorem C}
\newtheorem*{lemDP}{Lemma 1 (Dominated-Pair Lemma)}
\newtheorem*{lemJM}{Lemma 2 (join modes)}
\newtheorem*{connlem}{Connectivity lemma}
\newtheorem*{masterid}{Master identity}
\newtheorem*{finparts}{Finite-parts lemma}
\newtheorem*{advbij}{Advance bijection}
\newtheorem*{fullprof}{Full profile}

\newcommand{\rank}{\operatorname{rank}}
\newcommand{\elem}{\operatorname{elem}}
\newcommand{\cpc}{\operatorname{cap}}
\newcommand{\Z}{\mathbb{Z}}
\newcommand{\Q}{\mathbb{Q}}
\newcommand{\F}{\mathbb{F}}
\newcommand{\eO}{\mathcal{O}}
\newcommand{\PP}{\mathcal{P}}
\newcommand{\dd}{\delta}
\newcommand{\st}{\textsc}

\begin{document}

\title[Exact prime-count homology of the sphenic window complex]
{Exact prime-count homology of the sphenic window complex}
\author{Lee Rich}
\thanks{ORCID: \href{https://orcid.org/0009-0004-8559-5177}{0009-0004-8559-5177}.
The Improbability Lab. \emph{Working manuscript, 2 July 2026.}}
\email{leemjrich@gmail.com}
\date{2 July 2026}

\begin{abstract}
Fix a threshold $N$ and let the window be $W=(N/2,N]$. The \emph{sphenic window complex}
$K_N$ is the $2$-dimensional simplicial complex whose triangles are the sphenics (squarefree
integers $pqr$ with three distinct prime factors) lying in $W$, whose edges are the prime
pairs that co-occur in such a triangle, and whose vertices are the primes involved. These
complexes arise in connection with S.~Banerjee's question on the multiplicity of the
eigenvalue $-1$ of the integer coprime graph \cite{Banerjee2025}. We reduce the entire
homology of $K_N$ to prime-counting functions, with no linear algebra left in the answer, and
we determine an exact rational limiting density.

We prove a purely combinatorial identity (Theorem~A): for any finite \emph{slab} $R$, that is
any difference of two order ideals in the coordinatewise order on strictly increasing integer
tuples, the dimension of the top boundary-cycle space equals the number of \emph{fitting}
cells, those $x$ with the all-ones shift $x+(1,\dots,1)$ still in $R$, over every field, with
an explicit $\pm1$ integral basis. The sphenic complex is the length-$3$ slab in the
prime-position ruler, so Theorem~A gives $\rank\partial_2=F-\elem$ exactly (the rs-identity),
hence $b_1=(E-F)-V+1+\elem$. A finite-parts cancellation collapses this to a signed sum of
fourteen prime counts (Theorem~B), whose leading behaviour gives an exact rational density
$b_1/\pi(N/6)\to 6H-3/22=97360699/1078282205\approx0.0902924$ (Theorem~C), unconditional on
the prime number theorem, Bertrand's postulate, Dusart's explicit prime-gap bounds, and a
finite computation. The top Betti number is $b_2=\elem=F-D(N/2)+\cpc$ for every $N$, so the
full profile is $(b_0,b_1,b_2)=(1,b_1,\elem)$ for all $N\ge154$.

We close with an honest negative. The prime counts of Theorem~B carry the non-trivial zeros
of $\zeta$, and their amplitudes can be read off the fluctuation of $b_1$; this is exactly the
classical explicit-formula duality present in any prime count, and we verify directly that
only the prime ruler produces the zeros. We claim \textbf{nothing} toward the Riemann
hypothesis and no cheaper route to primes or zeros: the construction consumes the primes as
input.
\end{abstract}

\maketitle

\section{The sphenic window complex, and Banerjee's motivation}

\subsection{The object}
Fix $N$. The \emph{window} is $W=(N/2,N]$. A \emph{sphenic} is a squarefree integer with
exactly three distinct prime factors, $n=pqr$ with $p<q<r$. We build a $2$-complex $K_N$: the
\emph{vertices} are the primes dividing some sphenic in $W$ (to leading order, the primes up
to $N/6$); the \emph{triangles} ($2$-cells) are the sphenics in $W$, one $2$-simplex
$\{p,q,r\}$ each; the \emph{edges} are the prime pairs $\{p,q\}$ that co-occur in at least one
such triangle. Write $V,E,F$ for the numbers of vertices, edges and triangles, and let
$\partial_2$ be the simplicial boundary $[p,q,r]\mapsto[q,r]-[p,r]+[p,q]$. The two invariants
of interest are $\rank\partial_2$ over $\Q$ and the first Betti number
$b_1=E-(V-1)-\rank\partial_2$.

\subsection{Why this complex}
Banerjee \cite{Banerjee2025} studies the coprime graph of integers and gives a lower bound on
the multiplicity of the eigenvalue $-1$, leaving the exact value open. In the author's
Papers~1 and~2 \cite{Rich2026a,Rich2026b} that multiplicity is decomposed over the squarefree
layers of the integers, turning Banerjee's bound into an equality; the sphenic (three-prime)
layer contributes exactly the homological data of $K_N$, the first Betti number $b_1$ and a
companion count. The present paper, the third in that series, determines this layer: it pins,
exactly and rationally, the signed density underneath Banerjee's error term. It does not, on
its own, close the full multiplicity, which additionally requires the unsigned rank and the
other odd layers (see \S\ref{sec:shadow} and the companion Paper~4).

\subsection{The result in one paragraph}
The first Betti number is governed by a completely non-arithmetic fact about staircase-shaped
regions of a lattice (\S\ref{sec:slab}). Applied to the sphenic window that fact collapses,
after an exact cancellation, into a short signed sum of prime counts
(\S\ref{sec:slabsphenic}--\S\ref{sec:formula}). Dividing by the vertex count and letting $N$
grow gives an exact rational constant (\S\ref{sec:density}). The top Betti number is a single
semiprime count away from the same machinery (\S\ref{sec:top}). Finally
(\S\ref{sec:shadow}) we delimit the tempting appearance that the object ``sees'' the Riemann
zeros: it is the ordinary duality and nothing more.

\section{The combinatorial core: the slab identity}\label{sec:slab}

\subsection{Setup}
Fix an integer $\omega\ge2$. Let $P_\omega$ be the set of strictly increasing $\omega$-tuples
of positive integers, ordered coordinatewise: $x\le y$ iff $x_t\le y_t$ for every $t$. A
subset $D\subseteq P_\omega$ is a \emph{down-set} (order ideal) if $y\in D$ and $x\le y$ imply
$x\in D$. A \emph{slab} is a difference $R=D_{\mathrm{hi}}\setminus D_{\mathrm{lo}}$ of two
down-sets, $D_{\mathrm{lo}}\subseteq D_{\mathrm{hi}}$; every slab in this paper is finite.
Regard each $\omega$-tuple as an $(\omega-1)$-simplex on its entries, with the alternating
boundary $\partial=\partial_{\omega-1}$; restricting to the cells of $R$ gives a chain group
and a boundary map. Write $\dd=(1,\dots,1)$, call $x\in R$ \emph{fitting} if $x+\dd\in R$, and
write $x\ll y$ for strict domination, $x_t<y_t$ for every $t$.

\subsection{The theorem}
\begin{thmA}
\textnormal{\st{[proved; banked as C-0744]}}\ \ Fix $\omega\ge2$ and let
$R=D_{\mathrm{hi}}\setminus D_{\mathrm{lo}}$ be a finite slab in $P_\omega$. Then, over every
field,
\[
\dim\ker\partial|_R=\#\{x\in R: x+\dd\in R\}=(\text{number of fitting cells}).
\]
Moreover $\ker_{\Z}\partial|_R$ is free with an explicit basis of $\pm1$ vectors
(\S\ref{sec:assembly}), and every invariant factor of $\partial|_R$ equals $1$, so there is no
torsion anywhere in the Smith normal form.
\end{thmA}

\noindent The proof has two ingredients: a domination lemma giving the rank lower bound (the
non-fitting boundaries are independent), and an explicit family of cycles giving the matching
upper bound (each fitting cell carries a cycle expressible in strictly larger cells). Both are
field-agnostic and integral.

\subsection{Lemma 1: the Dominated-Pair Lemma}
\begin{lemDP}
\textnormal{\st{[proved]}}\ \ In the support of any nonzero $(\omega-1)$-cycle $z$, over any
field, there exist cells $x\ll y$. Equivalently, a $\ll$-antichain of cells has linearly
independent boundaries.
\end{lemDP}
\begin{proof}[Proof, base $\omega=2$]
A nonzero edge-chain with $\partial z=0$ gives every vertex of the support graph signed degree
zero, hence degree at least $2$, so the support contains a graph cycle $C$. Let $a^*$ be the
smallest vertex of $C$ and $(a^*,p),(a^*,q)$ its two edges on $C$, with $p<q$. Walk $C$ from
$p$ to $q$ avoiding $a^*$; every vertex on this walk exceeds $a^*$. Since the walk ends at
$q>p$ it contains an edge $g=(c,d)$ (sorted) with $d>p$ and $c>a^*$. Then $(a^*,p)\ll(c,d)$.
\end{proof}
\begin{proof}[Proof, induction step $\omega\ge3$ (assuming the lemma at $\omega-1$)]
Let $z\ne0$ with $\partial z=0$, and let $v$ be the smallest vertex in $\operatorname{supp}(z)$.
Split $\operatorname{supp}(z)=A\sqcup B$, where $A$ is the cells containing $v$ (as first
coordinate) and $B$ the rest, whose vertices all exceed $v$. Write each $x\in A$ as $v*x'$ with
tail $x'$ an increasing $(\omega-1)$-tuple on values above $v$; the tails are distinct. From
$\partial(v*x')=[x']-v*\partial[x']$, sort facets by whether they contain $v$:
\begin{itemize}
\item the $v$-part gives $v*(\sum_A c_x\,\partial x')=0$, and $v*$ is injective on $v$-free
chains, so $w:=\sum_A c_x[x']$ is a cycle, nonzero since $A\ne\varnothing$, the tails are
distinct, and the $c_x$ are nonzero;
\item the $v$-free part gives $\sum_A c_x[x']+\sum_B c_y\,\partial y=0$.
\end{itemize}
Apply the lemma at $\omega-1$ to $w$: there are tails $e\ll f$ in $\operatorname{supp}(w)$.
Since the coefficient of $[f]$ in $w$ is nonzero, the $v$-free relation forces $f$ to appear
in $\partial y$ for some $y\in B\cap\operatorname{supp}(z)$, say $y=f\cup\{s\}$ with $s\notin f$,
$s>v$.

\emph{Insertion claim: $(v,e)\ll y$.} The first coordinate satisfies $v<y_1$ since every entry
of $y$ exceeds $v$. Let $s$ land at sorted index $k$ of $y$, so $y_j=f_j$ for $j<k$, $y_k=s$,
$y_j=f_{j-1}$ for $j>k$. For each $t\in\{1,\dots,\omega-1\}$: if $t+1<k$ then
$y_{t+1}=f_{t+1}>f_t>e_t$; if $t+1=k$ then $y_k=s>f_{k-1}>e_{k-1}=e_t$; if $t+1>k$ then
$y_{t+1}=f_t>e_t$. In every case $e_t<y_{t+1}$, so $(v,e)\ll y$ with both cells in
$\operatorname{supp}(z)$.
\end{proof}
The step never uses the slab structure, and a degenerate base re-proves $\omega=2$, so the
whole lemma is one induction from a trivial base.

\subsection{Lemma 2: join modes}
\begin{lemJM}
\textnormal{\st{[proved]}}\ \ Let $x$ be a fitting cell of the slab $R$. Partition $x$ into
maximal runs of consecutive integers. For a run occupying values $a,a+1,\dots,b$, let its
factor be the boundary cycle of the simplex on $\{a,\dots,b+1\}$ (the alternating sum of its
subsets each omitting one value). Let $J_x$ be the join of the run factors, with product
signs. Then: \emph{(1)} $J_x$ is a nonzero cycle; \emph{(2)}
$\operatorname{supp}(J_x)=\{y: x\le y\le x+\dd\}$ exactly, and every such $y$ lies in $R$ by
the sandwich $y\le x+\dd\in D_{\mathrm{hi}}$ and $y\ge x\notin D_{\mathrm{lo}}$; \emph{(3)}
every coefficient of $J_x$ is $\pm1$, and $x$ is the unique coordinatewise minimum of
$\operatorname{supp}(J_x)$.
\end{lemJM}
Consequently, for any strictly monotone ruler (a value function that is a linear extension of
the coordinatewise order), every support cell of $J_x$ other than $x$ has strictly larger
value, so $\partial x$ lies in the span of $\{\partial y: y\in R,\ \mathrm{val}(y)>\mathrm{val}(x)\}$.
This resolves the apparent collision obstruction: a fitting cell with $x_{t+1}=x_t+1$ does not
carry a full $2^\omega$ box, but the degenerate box is a join of higher simplex boundaries (for
one collision, the suspension of a triangle boundary, a $2$-sphere on six cells).

\subsection{Assembly, and the integral basis}\label{sec:assembly}
\emph{Non-fitting cells form a $\ll$-antichain.} If $x,y\in R$ are non-fitting and $x\ll y$,
then $x+\dd\le y\in D_{\mathrm{hi}}$ and $x+\dd\ge x\notin D_{\mathrm{lo}}$ give $x+\dd\in R$, so
$x$ was fitting, a contradiction. \emph{$\rank\ge\#$non-fitting:} a nontrivial dependency among
the non-fitting boundaries is a nonzero cycle on a $\ll$-antichain, contradicting Lemma~1.
\emph{$\rank\le\#$non-fitting:} order the cells by strictly decreasing ruler value and
eliminate; when a fitting $x$ arrives, Lemma~2 writes $\partial x$ as a $\pm1$ combination of
strictly larger-value boundaries, already processed, so $x$ adds no pivot. Hence
$\rank=\#$non-fitting and $\dim\ker=\#$fitting, over every field.

For the integral statement, order the cells by any linear extension of the coordinatewise
order. For each fitting $x$, the vector $J_x$ has its least nonzero entry at row $x$ with value
$\pm1$ (Lemma~2), and distinct fitting cells give distinct pivot rows, so $\{J_x\}$ is in
permuted echelon form with unit pivots. This gives independence over every field; saturation of
the integer span (eliminate along the $\pm1$ pivots); and hence
$\ker_{\Z}\partial|_R=\Z\text{-span}\{J_x\}$ is free of rank $\#$fitting. Finally
$\rank_{\F_p}\partial=\#$non-fitting $=\rank_{\Q}\partial$ for every prime $p$, so no invariant
factor is divisible by any prime and the Smith normal form is $(1,\dots,1,0,\dots,0)$.

\subsection{Placement, and what is classical}
The strict-domination order is exactly the shifting order $x+\dd\le y$ of Klivans
\cite{Klivans}, underlying Kalai's algebraic shifting. \emph{The pure case is Björner--Kalai}
\cite{BjornerKalai}: for $D_{\mathrm{lo}}=\varnothing$ the complex is a near-cone, and their
theorem gives $\beta_{\omega-1}=\#\{\text{top cells avoiding vertex }1\}$; the bijection
$x\mapsto x-\dd$ carries these onto the fitting cells, so for pure down-sets Theorem~A is the
classical theorem in one line. \emph{The slab case is the new content, and it is not relative
homology:} the relative top cycles of $(D_{\mathrm{hi}},D_{\mathrm{lo}})$ need boundary only in
$D_{\mathrm{lo}}$, whereas our cycles need boundary zero in the ambient complex, a strictly
smaller space, so Duval's relative-shifting inequality \cite{DuvalRel} does not settle it. The
nearest machinery is Nagel--Reiner \cite{NagelReiner}: their Definition~3.4 names exactly our
slabs, and their complex-of-boxes cells are the same species as our join modes, but their
theorems compute the minimal free resolution of the associated ideal (a different functor);
their skew results are $d=2$ only, and the general-$d$ skew case is their open Question~5.1. The
honest verdict is that Theorem~A was not located there; the novelty is unclear pending a
specialist, and any external write-up must cite Björner--Kalai (pure case) and Nagel--Reiner
(box complexes). A residual check against Duval's relative Laplacian spectral recursion
\cite{DuvalLap} is the one remaining place the count could already be implicit.

\section{The sphenic complex as a slab, and the rs-identity}\label{sec:slabsphenic}

Number the primes $p_1=2,p_2=3,p_3=5,\dots$ and give each prime the coordinate equal to its
position $i$. A sphenic $p_ip_jp_k$ becomes the increasing triple $(i,j,k)$; its value
$p_ip_jp_k$ is used only to test window membership. The all-ones shift
$(i,j,k)\mapsto(i+1,j+1,k+1)$ is exactly ``advance each prime to the next prime,'' $p\mapsto
p^+$. Thus the sphenic complex is the $\omega=3$ slab in the prime-position ruler, and the
fitting condition is $p^+q^+r^+\le N$. Define the \emph{fitting count}
$\elem(N)=\#\{\text{sphenic }pqr\in W: p^+q^+r^+\le N\}$. Theorem~A at $\omega=3$ reads:
\[
\boxed{\ \rank\partial_2=F-\elem\ }\qquad\text{(the rs-identity), \st{[proved, corollary of Theorem~A]}.}
\]
Because $\rank\partial_2=F-\elem$, the first Betti number is a pure count:
$b_1:=E-(V-1)-\rank\partial_2=(E-F)-V+1+\elem$. That it equals the topological first Betti
number requires connectivity.
\begin{connlem}
\textnormal{\st{[proved; certified self-contained within the C-0743 banking-grade packet]}}\ \
For all $N\ge154$ the complex $K_N$ is connected, so $b_0=1$ and the count above is the Betti
number $b_1$ identically.
\end{connlem}
\noindent The proof is a hub argument through the vertex $2$: every prime appearing as a vertex
is tied to the $2$-hub through window sphenics guaranteed by Bertrand's postulate, Nagura's
bound, and Dusart's estimates once $N\ge N_c=154$, with the component count $c=1$ checked
exactly at four $N$. The separate edge-support layer-graph lemma (C-0739) is a related but
distinct result.

\section{An exact prime-count formula for $b_1$}\label{sec:formula}

\subsection{Ingredients}
Two facts are used as input; both sit inside the banked C-0743 package (independently
re-derived by hand, cleared through four hostile adjudication rounds, with a hashed audit
script and a formula-free brute-force confirmation at $N=1.05,1.20,1.50\times10^6$).
\begin{masterid}
\textnormal{\st{[proved; verified exact]}}\ \ With $B:=F-\elem$,
\[
E-B=\pi(N/4)+\cpc-D_E-1,
\]
where $\cpc$ counts the crossing triples $pqr\le N/2$ with $p^+q^+r^+>N$, and $D_E$ is a
blocker/isolation pair count. The transcendental semiprime term $Q_2(N/2)$ cancels identically,
which is why no Mertens-type constant survives.
\end{masterid}
\begin{finparts}
\textnormal{\st{[proved; certificate below]}}\ \ For $N>N_1=1{,}040{,}255$,
\[
\cpc=10+\!\!\sum_{(p,q)\in\mathcal T}\!\!\bigl[\pi(N/2pq)-\pi(N/p^+q^+)\bigr],\quad
D_E=\bigl[\pi(N/4)-\pi(N/6)\bigr]+\!\!\sum_{p\in\{3,5,7\}}\!\!\bigl[\pi(N/2p\,p^-)-\pi(N/p\,p^+)\bigr],
\]
where $\mathcal T=\{(2,3),(2,5),(2,7),(3,5),(3,7),(3,13),(3,19),(3,23),(5,7),(7,13)\}$ is the
set of ten crossing pairs.
\end{finparts}
\noindent The threshold $N_1$ comes from a finite family of at most $601$ late-crossing
triples, dominated by $(p^+/p)(q^+/q)\le1.98925<2$ together with Dusart's estimates, which
certify the largest such third prime satisfies $r<14{,}872$; the enumerated family's observed
maximum is $r=5{,}591$. The certified bound and the observed maximum are distinct numbers, both
kept on the record.

\subsection{The cancellation}
From $b_1=(E-F)-V+1+\elem=(E-B)-(V-1)$ and the master identity,
$b_1=\pi(N/4)-\pi(N/6)+\cpc-D_E$. Insert the finite-parts forms. The leading
$\pi(N/4)-\pi(N/6)$ cancels the $p=2$ term of $D_E$; the three cap-pairs $(2,3),(3,5),(5,7)$
cancel term for term against the blocker terms for $p=3,5,7$ (denominators $12/15$, $30/35$,
$70/77$ match exactly). Seven cap-pairs and the constant $10$ survive.
\begin{thmB}
\textnormal{\st{[proved, given the inputs above; verified]}}\ \ For all $N>1{,}040{,}255$,
\[
b_1(N)=10+\!\!\sum_{(p,q)\in\PP}\!\!\bigl[\pi(N/2pq)-\pi(N/p^+q^+)\bigr],\quad
\PP=\{(2,5),(2,7),(3,7),(3,13),(3,19),(3,23),(7,13)\}.
\]
\end{thmB}
\noindent Written out,
\begin{align*}
b_1(N)=10&+\pi(N/20)-\pi(N/21)+\pi(N/28)-\pi(N/33)+\pi(N/42)-\pi(N/55)\\
&+\pi(N/78)-\pi(N/85)+\pi(N/114)-\pi(N/115)+\pi(N/138)-\pi(N/145)\\
&+\pi(N/182)-\pi(N/187).
\end{align*}
The constant $10$ counts, for each crossing pair, the single completer dropped on the lower
edge of the open interval $\pi(N/2pq)-\pi(N/p^+q^+)$ (the largest prime $\le N/p^+q^+$, whose
successor clears $N$); these ten dropped completers are genuine cells, \st{[verified $10/10$ at
$N=2\times10^6$ and $5\times10^7$]}.

\subsection{Check}
The difference $b_1-(\text{seven-pair sum})$ equals $10$ exactly at nine values of $N$ from
$1.105\times10^6$ to $1.20\times10^8$, with $b_1$ computed independently as $(E-F)-V+1+\elem$
from the enumerated complex; the seven-pair sum reproduces the entire oscillation of $b_1$ to
numerical identity.

\section{The limiting density}\label{sec:density}

Each bracket satisfies $\pi(N/2pq)-\pi(N/p^+q^+)\sim(N/\log N)(1/2pq-1/p^+q^+)$, and
$V=\pi(N/6)\sim N/(6\log N)$.
\begin{thmC}
\textnormal{\st{[proved, unconditional on PNT $+$ Bertrand $+$ Dusart $+$ a finite
computation; banked as C-0743, banking-grade]}}
\[
b_1(N)/V(N)\ \longrightarrow\ K:=6H-\tfrac{3}{22}=\frac{97360699}{1078282205}=0.0902924\dots,
\]
where $H=\sum_{\text{ten pairs}}\bigl(1/2pq-1/p^+q^+\bigr)=\dfrac{488798363}{12939386460}$.
\end{thmC}
\noindent The limit is positive: writing $K=6\bigl(\tfrac14-\tfrac{7}{66}+H\bigr)-1$, positivity
is $H>1/44$, and $1/60+1/420+5/924$ already exceeds $1/44$. Equivalently, the three cancelled
pairs of $H$ sum to exactly $1/44$, so $K=6\times(\text{seven-pair sum})$, a sum of seven
positive brackets. No transcendental constant appears: the leading (Mertens-type) prime-count
constant cancels between the numerator's pieces, which is why $K$ is rational. Here $K$ is the
density of the signed first Betti number $b_1$; the lab separately tracks the canonical density
$(b_1-\dd)/V\approx0.0585$, where $\dd=\rank_{\text{uns}}\partial_2-\rank_{\text{sgn}}\partial_2$
is a torsion-like frustration term, not known to be rational and where any remaining
transcendence lives (the companion Paper~4 reduces $\dd$ to counting at a certified perimeter).
These two numbers must not be confused.

\section{The top Betti number, and the full profile}\label{sec:top}

The complex has no $3$-cells, so $b_2=\dim\ker\partial_2=\elem$ by Theorem~A directly. What
C-0745 adds is a closed form for $\elem$ through the advance map.
\begin{advbij}
\textnormal{\st{[proved; banked as C-0745]}}\ \ The next-prime map
$(p,q,r)\mapsto(p^+,q^+,r^+)$ carries increasing prime triples bijectively onto increasing
odd-prime triples (inverse: the previous-prime map, one-line surjectivity). The shell count
\[
J(N):=\#\{pqr\le N: p^+q^+r^+>N\}=D(N/2)=\#\{\text{odd semiprimes}\le N/2\}\quad\text{exactly, every }N.
\]
\end{advbij}
\noindent The lower-bound step is one line: $p^+q^+r^+>pqr>N/2$, so a shell triple's image is
an odd semiprime $\le N/2$, and the bijection matches the counts (recovering $168330$ at
$N=2\times10^6$). Splitting at the window edge gives, with no threshold,
\[
\rank\partial_2=D(N/2)-\cpc,\qquad b_2=\elem=F-D(N/2)+\cpc\quad\text{exact for every }N.
\]
Above the threshold this inherits the ten-pair $\pi$-form of $\cpc$: for $N>1{,}040{,}255$,
\[
b_2(N)=S(N)-S(N/2)-D(N/2)+10+\!\!\sum_{\text{ten pairs}}\!\!\bigl[\pi(N/2pq)-\pi(N/p^+q^+)\bigr],
\]
with $S$ the sphenic count, $F=S(N)-S(N/2)$. Verified exact at
$N\in\{2\times10^5,5\times10^5,1.05\times10^6,2\times10^6,5\times10^6,2\times10^7\}$; below $N_1$
the $\pi$-form is off by exactly $2$.
\begin{fullprof}
\textnormal{\st{[proved]}}\ \ For all $N\ge154$, $(b_0,b_1,b_2)=(1,b_1,\elem)$, with $b_1=10+$
fourteen prime counts for $N>1{,}040{,}255$, $b_1/V\to6H-3/22$, and
$b_2=\elem=F-D(N/2)+\cpc$. No linear algebra survives in any of the three numbers.
\end{fullprof}
\noindent For Banerjee's problem this supplies the three-prime layer in closed form: the signed
density his error term skirts is pinned exactly and rationally. The full $-1$ multiplicity still
needs the unsigned rank ($\dd$, addressed at the certified perimeter in the companion Paper~4) and
the other odd layers.

\section{The spectral shadow, and exactly how far it goes}\label{sec:shadow}

\subsection{What is true}
Because $b_1$ is a fixed signed combination of prime counts (Theorem~B), its fluctuation about
the smooth part is, by the classical explicit formula, a sum over the non-trivial zeros $\rho$
of $\zeta$:
\[
b_1(N)-(\text{smooth})=-\sum_\rho M(\rho)\cdot(N^\rho\text{-type term}),\qquad
M(\rho)=\sum_{m}\pm m^{-\rho},
\]
the sum over the fourteen denominators $m$. The factor $M(\rho)$, a Dirichlet polynomial in
fourteen small integers, is derived with no reference to $\zeta$. Scored against the measured
per-zero amplitudes of $b_1/V$ (data to $N=6\times10^{10}$, the known zeros used only as
sampling frequencies), it predicts the relative amplitudes with correlation $0.992$ over the
first $30$ zeros ($0.974$ without the dominant first zero), matching into the noise floor
($\gamma_{24}$ measured $0.026$ vs.\ predicted $0.027$; $\gamma_{30}$ $0.026$ vs.\ $0.028$),
\st{[verified]}. Reading zeros directly off the spectrum recovers $11$ of them, up to the
$18$th near $\gamma_{18}\approx72.1$, resolution-limited not signal-limited (the count grows
$6,9,11$ as sieve depth grows $3\times10^8,6\times10^9,6\times10^{10}$). The fluctuation grows
no faster than $N^{1/2}$ (measured exponent $0.44$), consistent with and required by the
Riemann hypothesis; the fourteen-term combination cancels about $99.6\%$ of the individual
prime-count fluctuation, so $b_1$ is an unusually quiet object.

\emph{The dartboard null} \st{[verified]}: tested against $20{,}000$ loose controls (fourteen
random distinct integers in $[10,200]$, balanced signs) and $20{,}000$ matched controls (seven
crossing-style pairs with ratio in $[1.005,1.6]$), the true envelope scores $0.998$ ($0.993$
excluding $\gamma_1$) against a null median $0.64$ and null maximum $0.94$ across all $40{,}000$
draws. With the shared decay stripped, the null collapses to median $\approx0.00$ while the true
envelope holds $0.992$. The amplitude channel is specific to the topology-selected denominators,
not an artifact of shared decay.

\subsection{What this is not}
All of the above is the ordinary prime-to-zero duality; the zeros are present because $b_1$ is
built from prime counts, and every prime count contains them (Riemann, 1859). Two controls make
the point. \emph{The envelope is blind to the zeros' positions:} $|M(\tfrac12+it)|$ averages
$0.489$ at the $30$ true zeros and $0.500$ at $5{,}000$ random heights, so the zeros sit at the
$48$th percentile of $M$'s own values. \emph{Only the prime ruler throws them:} the identical
construction with the plain-integer ruler peaks at the window-beat frequency, not the first
zero, while the prime ruler puts the first zero far above everything else. Consequently we claim
\st{[not claimed]} any progress on the Riemann hypothesis, any new bridge between primes and
zeros, or any cheaper route to either: computing $b_1$ requires the primes as input, and
nothing is obtained more cheaply than it is supplied.

\section{Discussion}

The self-contained core is the rational density $K=6H-3/22$ (Theorem~C) and its exact finite-$N$
form as fourteen prime counts (Theorem~B), resting only on PNT, Bertrand, Dusart, and a finite
certificate; this pins the signed density Banerjee's error term skirts, and the full profile
supplies the whole three-prime layer as counting functions. The most broadly interesting piece
may be Theorem~A itself, a statement with no arithmetic in it: its pure-down-set case is
classical (Björner--Kalai), and the slab case is a concrete theorem in a well-studied corner of
algebraic combinatorics; whether it is already on record is a question a specialist could settle
on sight, and the honest current label is ``adjacent, apparently unstated.'' What is honestly
not here is any route to the Riemann hypothesis or any predictor for primes.

\appendix
\section{The fourteen denominators}
Signed, in order: $+20,-21,+28,-33,+42,-55,+78,-85,+114,-115,+138,-145,+182,-187$. Each pair
$(2pq,p^+q^+)$ is a crossing pair with $p^+q^+>2pq$, so each bracket counts primes in a thin
interval and is non-negative.

\section{Certificates and data (reproducible)}
All entry points run from the repository root; the repository (concept
DOI \texttt{10.5281/zenodo.21135221}) holds the scripts and the certificate bundles under
\texttt{certificates/} as gzipped JSON with SHA-256 manifests.
\begin{itemize}
\item \emph{Slab identity and field-independence.} $\dim\ker$ over $\F_2,\F_3,\F_7,\F_{101},
\F_{2^{31}-1}$ agree (\texttt{slab\_gf2\_sweep.py}, $149$ slabs, $0$ mismatches); verified
$\omega=2,3,4$ across many rulers, $\omega=5,7$ for the prime case; proof verification
\texttt{dp\_proof\_verify.py} clean.
\item \emph{Theorem~B.} $b_1-$ seven-pair sum $=10$ at $N\in\{1.105,1.986,3.569,6.413,11.52,
20.71,37.21,66.87,120.2\}\times10^6$.
\item \emph{Theorem~C finite parts.} $H=488798363/12939386460$; the $D_E$ constant
$7/66=1/12+1/60+1/210+1/770$ with blockers $\{2,3,5,7\}$; $K=97360699/1078282205$; late-crossing
family $\le601$ triples; certified $r<14{,}872$ (Dusart) vs.\ observed maximum $r=5{,}591$;
hashed audit script reproduces all certificates plus a formula-free brute-force check at
$N=1.05,1.20,1.50\times10^6$.
\item \emph{Top Betti number.} \texttt{elem\_closed\_form\_verify.py}: $b_2=\elem=F-D(N/2)+\cpc$
exact at $N\in\{2\times10^5,5\times10^5,1.05\times10^6,2\times10^6,5\times10^6,2\times10^7\}$;
$J(N)=D(N/2)$ recovers $168330$ at $N=2\times10^6$.
\item \emph{Spectral shadow.} $30$-zero envelope correlation $0.992$ ($0.974$ without the first
zero); dartboard null over $40{,}000$ controls (\texttt{envelope\_dartboard.py});
position-blindness control $0.489$ vs.\ $0.500$; growth exponent $0.44$.
\end{itemize}

\section*{Acknowledgement / disclosure}
The reduction of the first Betti number to the slab fitting-count (the rs-identity) and the
finite-parts cancellation to fourteen prime counts are the author's. The manuscript was prepared
with assistance from a large language model (Claude, by Anthropic) for organisation, the
classical background and prior-art positioning, the slab formalism, and extensive symbolic and
numerical verification; all statements, proofs, and their interpretation are the author's and
were checked independently, with banking-grade claims filed through the laboratory's adversarial
adjudication protocol before being marked proved. Results carry status labels: \st{[proved]},
\st{[banked]}, \st{[verified]}, \st{[known]}, \st{[not claimed]}.

\begin{thebibliography}{9}
\bibitem{Rich2026a} L.~Rich, \emph{The $-1$-eigenspace of the coprime graph of integers:
prime-cube modes and an exact multiplicity}, Zenodo (2026), doi:10.5281/zenodo.20933665.
\bibitem{Rich2026b} L.~Rich, \emph{The $-1$-eigenvalue multiplicity of the coprime graph:
a conditional sublinear law}, Zenodo (2026), doi:10.5281/zenodo.21169867.
\bibitem{Banerjee2025} S.~Banerjee, \emph{On structural properties and adjacency spectrum of
coprime graph of integers}, Asian-Eur.\ J.\ Math.\ (2025), doi:10.1142/S179355712550069X;
arXiv:2506.10583.
\bibitem{BjornerKalai} A.~Björner, G.~Kalai, \emph{An extended Euler--Poincaré theorem}, Acta
Math.\ \textbf{161} (1988), 279--303.
\bibitem{Klivans} C.~Klivans, \emph{Threshold graphs, shifted complexes, and graphical
complexes}, Discrete Math.\ (2007), math/0703114.
\bibitem{DuvalRel} A.~M.~Duval, \emph{Algebraic shifting increases relative homology},
math/9809195, Discrete Math.\ (2000).
\bibitem{DuvalLap} A.~M.~Duval, \emph{A relative Laplacian spectral recursion}, Electron.\ J.\
Combin.\ \textbf{11}(2) (2006), \#R26, math/0507130.
\bibitem{NagelReiner} U.~Nagel, V.~Reiner, \emph{Betti numbers of monomial ideals and shifted
skew shapes}, Electron.\ J.\ Combin.\ \textbf{16}(2) (2009), \#R3, arXiv:0712.2537.
\bibitem{Dusart} P.~Dusart, \emph{Explicit estimates of some functions over primes}, Ramanujan
J.\ \textbf{45} (2018), 227--251.
\end{thebibliography}

\end{document}
